\section{Combining the arithmetic and real estimates}
\label{s7-completion}

This section is the analogue of source Section~7, p.~21.  Its task is to
put the two quadratic coefficients on the same scale.  We retain
$K=40n$, $N=3n$, $h=37n$, and write $\xi=\zeta(9)$.  The outcome is a
positive upper-bound budget; the final small-value step of the source
therefore becomes a conditional implication.

\subsection{The normalized polynomial and its budget}
\label{s7-budget}

\noindent\textbf{[Conditional]} With the arithmetic estimates used in
Section~5, let
\begin{equation}
 Q_{9,K,M}(X)=m_{9,K,M}F_{9,K}(X)
            =m_{9,K,M}S_K\Delta_{9,K}(X)\in\mathbb Z[X],
 \qquad m_{9,K,M}>0.
 \label{s7-normalized}
\end{equation}
The degree is $h$ by the leading-coefficient calculation in Section~2.
The positive Gram representation and the positive scalar factors give
$Q_{9,K,M}(\xi)>0$.  The normalization changes neither fact.  For fixed
$M\in40\mathbb Z_{>0}$ the input estimates are
\[
 \log m_{9,K,M}\leq A_{9,M}K^2+o_M(K^2),\qquad
 \log F_{9,K}(\xi)\leq UK^2+28K\log K+200K,
 \qquad U=-\frac{2733991}{2000000}.
\]
The second estimate is \eqref{s6-final}; it uses the source Appendix~A
numerical certificate discussed in Appendix~\ref{appa-measure}.
Adding, and then dividing by $K^2$, gives the
direct analogue of source (7.1):
\begin{equation}
 \limsup_{\substack{K\to\infty\\40\mid K}}
  \frac{\log Q_{9,K,M}(\xi)}{K^2}
 \leq B_{9,M},\qquad B_{9,M}:=A_{9,M}+U.
 \label{s7-combined}
\end{equation}
Thus the arithmetic loss has to be paid before the negative real
coefficient can be used to obtain a small integer polynomial value.
The order of the limits matters here: $M$ is held fixed while $K$
increases through multiples of forty.  The comparisons between cutoffs
below compare the resulting constants.  They do not require a uniform
error estimate with a cutoff $M$ depending on $K$.  In particular, the
explicit $K\log K$ error in the real estimate does not upgrade the
arithmetic remainder $o_M(K^2)$ to an error of that same order.

\noindent\textbf{[Derived]} The arithmetic coefficient obtained with the
unchanged $D_N^6$ construction is
\begin{align}
 A_{9,M}&=A_{5,M}+\frac{116581}{57600}
 \notag\\
 &=A_{9,*}+\frac{7\lambda}{M}
       -\frac{2923/240-1/4}{M^2}+\frac{32}{M^3},
 \qquad \lambda=\frac{37}{40},
 \label{s7-arithmetic-coefficient}\\
 A_{9,*}&=\frac{25180312106351244019069}
                   {7535670527041937280000}
           \simeq3.341482621353.
 \label{s7-limit-coefficient}
\end{align}
Here $A_{5,M}$ denotes the source coefficient, rather than a new
fifth-order construction.  The rational shift summarizes the increased
denominator cost calculated in Section~5.  In particular, the comparison
replacing the negative margins of source (7.2) is
\begin{equation}
 \begin{array}{c|c|c}
 M&A_{9,M}&B_{9,M}=A_{9,M}+U\\ \hline
 200&3.373563392186\ldots&2.006567892186\ldots\\
 100000&3.341547370160\ldots&1.974551870160\ldots\\
 \infty&3.341482621353\ldots&1.974487121353\ldots
 \end{array}
 \label{s7-sign-table}
\end{equation}
Increasing the cutoff does not change this sign.  Indeed, the correction
to $A_{9,*}$ is
\[
 \frac{1554M^2-2863M+7680}{240M^3}>0\qquad(M\geq40),
\]
and even $A_{9,*}+U$ is positive.

\subsection{The conditional final step}
\label{s7-conditional-step}

\noindent\textbf{[Blocked conclusion]} Substitution of $K=40n$ in
\eqref{s7-combined} supplies an upper coefficient $1600B_{9,M}>0$ on
the $n^2$ scale.  It supplies no decay rate of the type used in source
(2.7)--(2.8).  This does not prove that the polynomial values grow: a
positive upper bound alone does not decide their actual size.  It says
precisely that the present arithmetic and real estimates do not close
the small-value argument.

\noindent\textbf{[Conditional]} The remaining manipulation can still be
recorded.  Suppose some strengthened estimates, for a fixed cutoff,
gave an $\varepsilon>0$ and
\begin{equation}
 0<Q_{9,K,M}(\xi)\leq e^{-\varepsilon K^2}
 \quad\text{for all sufficiently large }K\in40\mathbb Z_{>0}.
 \label{s7-hypothetical-decay}
\end{equation}
If $\xi=a/b$ with $a\in\mathbb Z$ and $b\in\mathbb Z_{>0}$, degree
$h=\lambda K$ and integrality would imply
\[
 0<b^hQ_{9,K,M}(a/b)
 \leq\exp\{\lambda K\log b-\varepsilon K^2\}\longrightarrow0,
 \qquad b^hQ_{9,K,M}(a/b)\in\mathbb Z.
\]
This is the source's final contradiction, with the moment order changed
to nine.  Its algebra is unchanged, but its negative-budget hypothesis
is not furnished by \eqref{s7-sign-table}.  Appendix~C likewise records
the rational-approximation implication under an explicit
$\varepsilon$ hypothesis, while retaining the relative norm and height
calculations that can be made with the present budget.
